\section{Introduction}\label{sec:introduction}

% P1
Let \(n\geq8\) be even.  We write
\[
G_n=C_n^2=C_n(1,2)
\]
for the graph on \(\mathbb Z/n\mathbb Z\) in which two vertices are adjacent
when their cyclic distance is one or two.  An edge signing is a map
\(\sigma:E(G_n)\to\{\pm1\}\), and \(A_\sigma\) is the real symmetric matrix
whose \(i,j\) entry is \(\sigma(ij)\) on an edge and zero otherwise.  Our
problem is to determine
\[
m_n=\min_{\sigma:E(G_n)\to\{\pm1\}}\rho(A_\sigma).
\]
Thus the already formed cycle square is being signed: the distance-one and
distance-two edges receive signs independently.  It is not the graph obtained
by first signing a cycle and then squaring that signed object.  Switching at
vertices conjugates \(A_\sigma\) by a diagonal \(\{\pm1\}\)-matrix, so the
minimization naturally takes place over switching classes.

% P2
The problem belongs to the general signing program in spectral graph theory.
Signings encode the new eigenvalues of two-lifts in the work of Bilu and
Linial \cite{BiluLinial2006}, while Belardo, Cioab\u{a}, Koolen, and Wang ask
explicitly which signatures minimize the adjacency spectral radius of a fixed
connected graph \cite{BelardoCioabaKoolenWang2019}.  Interlacing families give
strong existence results for one side of the signed spectrum; in the
bipartite setting spectral symmetry turns this into the two-sided control
used to construct Ramanujan lifts \cite{MarcusSpielmanSrivastava2015}.  The
graphs \(C_n^2\) are nonbipartite, however, so a one-sided eigenvalue estimate
does not by itself control \(\rho(A_\sigma)\).  Here the graph is fixed at each
order and the question is the exact two-sided optimum, rather than an
existence bound uniform over a broad graph class.

% P3
The problem-specific starting point is Suvagiya's preprint
\cite[Conjecture~3]{Suvagiya2026Signed}.  For \(C_n(1,2)\), it identifies four distinguished
switching classes in which every quadrilateral is unbalanced, computes the
spectra of the two twisted classes, verifies global optimality for even orders
through \(18\), and formulates the all-even-order assertion as Conjecture~3.
We determine exactly when that twisted candidate is globally minimizing,
thereby resolving and disproving the conjecture.  In particular, the problem,
the flux coordinates, and the twisted spectral formula are inputs from that
work, not novelty claims of the present paper.

% P4
The answer is unexpectedly nonmonotone.  The proposed equality holds for
every even order from \(8\) through \(30\), fails first at \(32\), returns at
\(34,36,38\), fails again at \(40\), and returns once more at \(42,44,46\).
Only from \(48\) onward does failure persist at every even order.  The first
counterexample therefore does not mark the beginning of continuous failure,
and the two isolated failures cannot be inferred from the eventual
phase-slip mechanism developed below.

% P5
To state the classification, set
\[
\rho_-(n)^2
=4+2\cos\frac{2\pi}{n}+2\cos\frac{4\pi}{n},
\qquad
\theta_n=\rho_-(n)^2.
\]
The distinguished class has alternating triangle flux and negative
Hamilton-cycle holonomy; a direct antiperiodic Fourier calculation shows that
it attains \(\rho_-(n)\) for every even \(n\geq8\).  This gives the upper
bound \(m_n\leq\rho_-(n)\) at all orders.  At an order where equality holds,
the reverse inequality for every signing is a separate conclusion of exact
exhaustion.

\begin{theorem}[Complete classification]\label{thm:complete-classification}
For every even \(n\geq8\), the distinguished twisted signing attains
\(\rho_-(n)\), and
\[
m_n<\rho_-(n)
\quad\Longleftrightarrow\quad
n=32,\quad n=40,\quad\text{or}\quad n\geq48.
\]
Equivalently,
\[
m_n=\rho_-(n)
\quad\Longleftrightarrow\quad
n\in\{8,10,12,14,16,18,20,22,24,26,28,30,34,36,38,42,44,46\}.
\]
\end{theorem}

Theorem~\ref{thm:complete-classification} classifies the truth of the
proposed equality.  It neither classifies all minimizing signings at the
valid orders nor determines the exact value of \(m_n\) at the failing orders.

% P6
Its finite component determines the irregular small-order truth table.  A
different part of the argument explains why failure is eventually
uninterrupted: a low-spectral periodic phase supports localized elementary
interfaces, and well-separated copies of those interfaces can be closed on
large rings.  The assertion that continuous failure begins at \(48\) uses
both this analytic tail and an exact finite bridge; the interface theory
alone gives no such sharp onset.

% P7
The structural comparison phase is period eight and is distinct from the
twisted signing that attains \(\rho_-(n)\).  In its quadrilateral-flux word,
the positive sites occur at spacing four.  Replacing one spacing by a
positive gap \(g\) creates a single interface; we call \(q=g-4\) its excess
charge.  For a cyclic collection of gaps, the global closing law is
\(\sum_j q_j\equiv n\pmod 8\), whereas the local translation sector carried
by one charge is \(\sigma_{\mathrm{sec}}(q)=q\pmod4\).  These congruences have
different roles and are not interchangeable.  In particular, one, two, and
three copies of the gap-six interface, whose charge is two, supply the three
nonzero even residue classes needed in the finite-ring constructions.

% P8
The reference squared spectral edge is
\[
\eta=4+\sqrt{10+2\sqrt5}.
\]
For the bilateral positive single-gap operator \(A_g\), put \(H_g=A_g^2\)
and \(E(g)=\sup\operatorname{Spec}(H_g)\).  Define \(c_6\) to be the unique
root in the rational interval
\[
\frac{7905369311620327}{10^{15}}<c_6<
\frac{7905369311620328}{10^{15}}
\]
of the degree-ten polynomial
\[
\begin{aligned}
p_6(y)={}&16y^{10}-520y^9+6913y^8-48448y^7+191768y^6\\
&-423904y^5+484528y^4-270464y^3+137856y^2\\
&-19968y+256.
\end{aligned}
\]
The exact isolation is important: all later strict comparisons are made
against rational endpoints, not against a decimal approximation.

\begin{theorem}[Reference and single-gap spectral hierarchy]
\label{thm:single-gap-hierarchy}
For both lifts and both orientations,
\[
E(4)=\eta<c_6=E(6),
\qquad
\operatorname{Spec}_{\mathrm{ess}}(H_6)
=\operatorname{Spec}(H_{\mathrm{ref}}),
\qquad
\dim\ker(H_6-c_6)=2.
\]
Moreover, for every positive integer \(g\notin\{4,6\}\),
\[
E(g)>c_6+\frac1{250}.
\]
\end{theorem}

% P9
The scope of Theorem~\ref{thm:single-gap-hierarchy} is precise: G6 is the
unique minimizer among abnormal positive single gaps.  It is not asserted to
minimize over arbitrary multi-gap or finite-core interfaces.  The strict
single-gap hierarchy makes G6 the elementary low-cost defect above the
period-eight bulk.  Its charge allows one, two, or three separated copies to
close in the required residue classes.  Exact patch identification transfers
the bilateral reference/G6 bounds to those finite rings, and a discrete IMS
partition controls the squared spectral top between the patches.  This proves
failure for every even \(n\geq240\).  Exact finite verification on
\(48\leq n<240\) then shows that continuous failure begins precisely at
\(48\).

% P10
The computer-assisted part is organized as a proof by finite reduction, not
as numerical evidence.  Switching quotients and local interlacing reduce the
relevant cases to finite exact objects.  For \(8\leq n\leq32\), complete
switching-class decisions establish the lower bounds through \(30\) and an
exact witness at \(32\).  For \(34,36,38,42,44,46\), a parity-lifted finite
state graph is proved sound, complete, and terminating, and every terminal
case is closed exactly; a separate rational positive-definiteness certificate
gives the witness at \(40\).  Finally, rational positive-definiteness
certificates cover every even order from \(48\) through \(238\).  Independent
checkers reconstruct the finite objects and repeat the integer, rational, or
algebraic decisions.  Exhaustive generation and local finite-state reductions
have established precedents in graph theory \cite{GoedgebeurSchaudt2018,DeVosSamal2011},
while all-order classifications often combine an analytic tail with finite
closure \cite{LinNing2021,GoedgebeurRendersWienerZamfirescu2024}.

% P11
Exact signed spectral-radius classifications are known in settings where the
underlying graph is also allowed to vary, including unbalanced unicyclic and
bicyclic graphs \cite{BelardoBrunettiCiampella2021} and connected unbalanced
signed graphs of fixed order \cite{BrunettiStanic2022}.  Those results make it
essential to keep the present novelty claim narrow.  To the best of our
knowledge, no previous work gives an all-order classification of the
minimizing behavior over all real edge signings of the fixed circulant family
\(C_n(1,2)\).  The closest direct source remains Suvagiya's conjecture; the
complete truth set and the elementary single-gap mechanism are the two
contributions supplied here.  The paper is arranged as follows.
Section~\ref{sec:switching-reference}
sets up switching coordinates, the twisted attaining candidate, and the
period-eight reference phase.  Section~\ref{sec:gaps-charges-sectors}
develops gaps, charges, and the two closing laws.  Sections~\ref{sec:elementary-g6}
and \ref{sec:single-gap} prove the G6 edge theorem and the complete
single-gap hierarchy.  Section~\ref{sec:finite-rings} passes from bilateral
interfaces to finite rings by patch identification and discrete IMS
localization.  Section~\ref{sec:finite-completion} closes the remaining
orders by exact finite arguments, and Section~\ref{sec:concluding} records
the resulting perspective.  The appendices collect the algebraic G6
certificate and the exact finite-classification boundary.
